%% Lista de Pacotes utilizados

                                             %% Recomendacao: 11pt, papel a4
\usepackage{ifthen}                          %% Nao remover. Uso interno do sistema.
\usepackage[T1]{fontenc}                     %% Acentua��o
%\usepackage[portugues]{babel}
%\usepackage[latin1]{inputenc}               %% Use um dos dois esquemas

\usepackage[cp1252]{inputenc}                %% Use um dos dois esquemas


\usepackage{type1cm}                         %% Liberdade para as fontes!!!
\usepackage{lastpage}                        %% Refer�nciar a �ltima p�gina
\usepackage{textcomp}
\usepackage{nfssext}
\usepackage[labelfont={bf},textfont={it}]{caption}     %% Controlar melhor os captions....
                                                       %% Deixando o caption mais interessante...

\usepackage{times}                          %% Se estiver achando muito bonito,
                                             %% e quizer piorar um pouco,
                                             %% usar fonte Times
%\usepackage[oldstyle]{sabon}                 %%
%\usepackage{mathpazo}                        %% Fonte parecida com a sabon !
%\usepackage{eco}                             %% N�meros em oldstyle em Computer Modern

\usepackage{amsmath}                         %% [leqno]
\usepackage{amsfonts,bm}                     %% Pacotes da AMS
\usepackage{amssymb}                         %%
\usepackage{indentfirst}
\usepackage{enumitem}
\usepackage{array,arydshln,colortbl}
\usepackage{lscape}
%\usepackage[table]{xcolor}
%\usepackage[space]{cite}
\usepackage[square,comma,numbers,sort&compress]{natbib} %%
%\usepackage{longtable}                       %% Para a lista de abreviaturas
\usepackage{wrapfig}
\usepackage[english, portuges]{babel}         %% Hifenizacao e estilo
\usepackage{float}                           %% Descomente caso precise de ``floats'' extra
                                             %% ex: fragmentos de programas
\usepackage{graphicx}                        %% Mais ferramentas para
\usepackage{psfrag}                          %% manuseio de figuras.
\usepackage{subfigure}
%\usepackage{caption}
%\usepackage{subcaption}

\usepackage{colordvi}                        %% Usar Cores...
\usepackage{fancyhdr}
\usepackage{yfonts}                          %% n�o mais restri��es de tamanhho!
\usepackage{pstricks}                        %% Pacote de comandos Postscripts
\usepackage{pstricks-add}                    %% " "
\usepackage{pst-all,pst-math,pst-plot}       %% " "

%\usepackage{draftwatermark}                 %% Caso queira inserir uma marca d'agua

\usepackage[T1]{pbsi}

%\usepackage{enumerate}                       %% Pra enumera��o de t�picos...
%\usepackage{mathrsfs}                        %% $\mathscr{ABCD}$
                                             %% Pra certos s�mbolos ficarem melhores...
                                             %% como o de L^2...
\usepackage[amsmath,thmmarks]{ntheorem}      %% Facilita a personaliza��o do ambiente para teoremas e fornece mais op��es

\usepackage{multirow}                        %% Juntar linhas em tabelas
\usepackage{eucal}                           %% beau
\usepackage{lettrine}                        %% Pra colocar a primeira letra do cap�tulo
                                             %% maior.
\usepackage{footmisc}

\usepackage{url}                             %% Para URL, PATH, etc
\usepackage{mathrsfs}
\usepackage{algorithm}
\usepackage{varwidth}
\usepackage[noend]{algpseudocode}
\usepackage{amssymb}
\usepackage{tikz}
\usetikzlibrary{graphs, graphs.standard}

%\newdimen\pgf@lib@svg@lastmoveto@x
%\newdimen\pgf@lib@svg@lastmoveto@y
%
%\def\pgf@lib@svg@moveto{%
%  \pgf@lib@svg@clear@bezier@quad%
%  \pgf@lib@svg@last@x=\pgf@lib@svg@get@num{0}pt%
%  \pgf@lib@svg@last@y=\pgf@lib@svg@get@num{1}pt%
%  \pgf@lib@svg@lastmoveto@x=\pgf@lib@svg@last@x% NEW LINE
%  \pgf@lib@svg@lastmoveto@y=\pgf@lib@svg@last@y% NEW LINE
%  \pgfpathmoveto{\pgfqpoint{\pgf@lib@svg@last@x}{\pgf@lib@svg@last@y}}%
%  \pgf@lib@svg@read@nums{2}{\pgf@lib@svg@lineto}
%}
%\def\pgf@lib@svg@moveto@rel{%
%  \pgf@lib@svg@clear@bezier@quad%
%  \advance\pgf@lib@svg@last@x by\pgf@lib@svg@get@num{0}pt%
%  \advance\pgf@lib@svg@last@y by\pgf@lib@svg@get@num{1}pt%
%  \pgf@lib@svg@lastmoveto@x=\pgf@lib@svg@last@x% NEW LINE
%  \pgf@lib@svg@lastmoveto@y=\pgf@lib@svg@last@y% NEW LINE
%  \pgfpathmoveto{\pgfqpoint{\pgf@lib@svg@last@x}{\pgf@lib@svg@last@y}}%
%  \pgf@lib@svg@read@nums{2}{\pgf@lib@svg@lineto@rel}
%}
%
%\def\pgf@lib@svg@closepath{
%  \pgf@lib@svg@finish@prev
%  \pgf@lib@svg@last@x=\pgf@lib@svg@lastmoveto@x% NEW LINE
%  \pgf@lib@svg@last@y=\pgf@lib@svg@lastmoveto@y% NEW LINE
%  \pgfpathclose
%  \let\pgf@lib@svg@finish@prev=\relax
%  \pgf@lib@svg@clear@bezier@quad%
%  \pgfparserswitch{initial}%
%}

\usetikzlibrary{trees}
\usetikzlibrary{positioning}
%\usepackage{euler}
%\usepackage[dvips]{hyperref}

%\renewcommand{\captionfont}{\bfseries}
\addto\captionsenglish{\renewcommand{\chaptername}{Chapter}}
\addto\captionsenglish{\renewcommand{\figurename}{Figure}}
\addto\captionsenglish{\renewcommand{\tablename}{Table}}



\hfuzz2pt
%
%\hypersetup{backref,%
%            pdfpagemode=FullScreen,%
%            colorlinks=true,%
%            citecolor=orange,%
%            urlcolor=blue, %
%            pdfauthor={Jos� Sampaio de Lemos Neto}, %
%            pdftitle={Constru��o de Sequ�ncias de Protocolo para o Canal de Colis�o sem Realimenta��o}, %
%            pdfcreator={Jos� Sampaio de Lemos Neto}, %
%           breaklinks=true, %
%           linktocpage=true
%            }%

%%%%%%%%%%% Estilos definidos por Jos� Sampaio de Lemos Neto em 26/12/2010 %%%%%%%%%%%%%%%%%%%%%%%%%%%%%%

\makeatletter

%% Meu estilo para Teoremas, Defini��es, Exemplos, etc. -----------------------------------------------
\newtheoremstyle{sampaio}
{\vspace{0.4cm}\item[\rlap{\hspace{-0.6cm}\vbox{\hbox{\hskip\labelsep \theorem@headerfont
##1\ ##2\theorem@separator}\hbox{\strut}}}]}
{\vspace{0.4cm}\item[\rlap{\hspace{-0.6cm}\vbox{\hbox{\hskip\labelsep \theorem@headerfont
##1\ ##2\ -- \textmd{##3}\theorem@separator}\hbox{\strut}}}]}
%% ----------------------------------------------------------------------------------------------------

%% Estilo de refer�ncias utilizado pelo IEEEtrans.cls. Modifica��es executadas sobre o pacotes cite.sty
\def\@citex[#1]#2{%
  \let\@citea\@empty
  \@cite{\@for\@citeb:=#2\do
    {\@citea\def\@citea{], [}%
     \edef\@citeb{\expandafter\@firstofone\@citeb\@empty}%
     \if@filesw\immediate\write\@auxout{\string\citation{\@citeb}}\fi
     \@ifundefined{b@\@citeb}{\mbox{\reset@font\bfseries ?}%
       \G@refundefinedtrue
       \@latex@warning
         {Citation `\@citeb' on page \thepage \space undefined}}%
       {\hbox{\csname b@\@citeb\endcsname}}}}{#1}}
\def\citepunct{], [}
\def\citedash{]--[}
%% ---------------------------------------------------------------------------------------------------

\makeatother

\theoremstyle{sampaio}
\theoremindent0.6cm
\theoremsymbol{$\square$}
\theoremheaderfont{\normalfont\bfseries}
\theorembodyfont{\it}

\newtheorem{proposition}{Proposition}[chapter]  %%
\newtheorem{definition}{Definition}[chapter]    %%
\newtheorem{property}{Property}[chapter]            %%
\newtheorem{theorem}{Theorem}[chapter]                 %%
\newtheorem{lemma}{Lemma}[chapter]                      %%
\newtheorem{corollary}{Corollary}[chapter]         %%
\newtheorem{conjecture}{Conjecture}[chapter]           %%

\newcommand{\etal}{\emph{et\:al.}}
\newcommand{\fonte}{\footnotesize Source: Author.}

\theoremstyle{sampaio}
\theoremindent0.6cm
\theoremsymbol{$\square$}
\theoremheaderfont{\normalfont\bfseries}
\theorembodyfont{\it}

\newtheorem{example}{Example}[chapter]                 %%
\newtheorem{construction}{Construction}[chapter] %%
\newtheorem{procedure}{Procedure}[chapter]           %%

\theoremstyle{nonumberplain}
\theoremheaderfont{\it}\theorembodyfont{\normalfont}
\theoremseparator{:}
\theoremsymbol{$\blacksquare$}
\theoremindent0.6cm

\newtheorem{proof}{\hspace{-0.85cm} Proof}%%
\theoremstyle{nonumberbreak}
\theoremheaderfont{\it}\theorembodyfont{\normalfont}
\theoremseparator{:}
\theoremsymbol{$\blacksquare$}
\theoremindent0.6cm

\newtheorem{proof2}{\hspace{-.7cm} \underline{Proof}}%%
\theoremheaderfont{\sc}
\theorembodyfont{\upshape}
\theoremstyle{nonumberplain}
\theoremseparator{.}
\theoremsymbol{$\blacksquare$}
%\theoremindent=-0.4cm



%% Para personalizar as tabelas

\newcolumntype{I}{!{\vrule width 1pt}}
\newcolumntype{|}{!{\vrule width 0.5pt}}
\newlength\savedwidth
\newcommand\whline{\noalign{\global\savedwidth\arrayrulewidth
\global\arrayrulewidth 1.5pt}%
\hline
\noalign{\global\arrayrulewidth\savedwidth}}
\setlength\dashlinedash{0.25pt}
\setlength\dashlinegap{1.5pt}

%% Personalizar a Marca d'agua

%\SetWatermarkLightness{0.9}
%\SetWatermarkFontSize{2cm}
%\SetWatermarkScale{1.2}
%\SetWatermarkText{Vers�o Preliminar}

%%%%%%%%%%%%%%%%%%%%%%%%%% FIM DAS MINHAS DEFINI��ES (Sampaio) %%%%%%%%%%%%%%%%%%%%%%%%%%%%%%%%%%%%%%%%%%%%%%%%%%%%%%%%%%%%%%%%%

\numberwithin{equation}{chapter}                       %%
\numberwithin{table}{chapter}                          %% Numera��o dos floats...
\numberwithin{figure}{chapter}                         %%
\DeclareMathAlphabet{\mathds}{U}{eus}{m}{n}            %% Definindo o comando \mathds... (Jos� Sampaio de Lemos Neto em 29/12/2010)
\DeclareMathAlphabet{\mathcal}{OMS}{cmsy}{m}{n}        %% Para o mathcal ficar bom... (Jos� Sampaio de Lemos Neto em 29/12/2010)
\DeclareMathAlphabet\mathbb{U}{fplmbb}{m}{n}           %% Para o mathbb ficar bom...

\DeclareMathOperator*{\sinc}{sinc}                     %% Declarando novos operadores
\DeclareMathOperator*{\sen}{sen}                       %%
\DeclareMathOperator*{\cas}{cas}                       %%

\setlist{noitemsep}
\setlist{labelindent=\parindent} % << Usually a good idea
\setlist{itemsep=0pt}%
\setitemize{leftmargin=*} %
\setitemize[1]{label=$\triangleright$} %
\setenumerate[1]{label=\arabic*., ref=\arabic*}
%\setenumerate[1]{label=\Roman*., ref=\arabic*}
\setenumerate[2]{label=\emph{\alph*}), ref=\theenumi.\emph{\alph*}}
\setenumerate[3]{label=\roman*), ref=\theenumii.\roman*}
%\setdescription{font=\sffamily\bfseries}

%%% Criando novas funcoes
\newcommand{\unidade}[2]{\ensuremath{#1\,\mathrm{#2}}}

% Caso a tese tenha fragmentos de c�digo no corpo da disserta��o,
% descomente as linhas abaixo e use o ambiente
% \begin{programa} ... \end{programa}
%%%%% Inicio da defini��o de ambiente para programas
\floatstyle{ruled}
\newfloat{Program}{thp}{lop}[chapter]
\floatname{Program}{Programa}

\newenvironment{programa}[1]
{
\begin{Program}
\caption{#1} } {\end{Program}}

%%%%% Fim da defini��o de ambiente para programas %%%%%%%

%\def\QED{\mbox{\rule[0pt]{1.5ex}{1.5ex}}\vspace{2mm}}
%\def\proof{\hspace{0.6cm}{\it Demonstra\c{c}\~{a}o: }}
%\def\endproof{\hspace*{\fill}~\QED\par\endtrivlist\unskip}


\renewcommand{\headrulewidth}{0pt}%
\renewcommand{\figurename}{Figure}

\newsavebox{\mygraphic}
%\sbox{\mygraphic}{\includegraphics[scale=1.2]{Figuras/ufpe.eps}}


\newenvironment{my-notation}{
    \begin{list}{}{\setlength{\labelwidth}{55pt}
                   \setlength{\labelsep}{40pt}
                   \setlength{\leftmargin}{15em}}}
    {\end{list}}
